\section{AI Accelerators - Dataflow}\label{chap:background:sec:dflow}

This section describes the second class of Domain-Specific Accelerators for LLM inference, which consists of \textbf{dataflow accelerators}. \ref{tab:dflow} summarizes the main features of the architectures discussed in this section, namely SambaNova RDAs, Cerebras WSE, and GraphCore IPUs.

\subsection{SambaNova RDA}
\label{chap:background:sec:sambanova}

Founded in 2017, \textbf{SambaNova Systems} develops dataflow accelerators for DNN training and inference, inspired by the \textbf{Plasticine} prototype architecture~\cite{Prabhakar2017}.

\textbf{Plasticine} is a coarse-grained reconfigurable architecture (CGRA) that exploits spatial parallelism for compute patterns such as map and reduce. Unlike bit-level FPGAs, CGRAs use word-level processing elements (ALUs, DSPs) connected through structured interconnects, achieving higher density and clock speeds. Plasticine consists of \textbf{Pattern Compute Units (PCUs)} and \textbf{Pattern Memory Units (PMUs)} arranged in a grid, managed by hierarchical controllers that coordinate pipelined and streaming execution. Address Generators and Coalescing Units handle off-chip memory, optimizing access patterns and reducing traffic.

Building on this concept, SambaNova’s \textbf{Reconfigurable Dataflow Architecture} \cite{sambanova} combines tiled PCUs and PMUs linked by a high-bandwidth switch fabric, with AGUs and CUs managing external memory. Its software stack, \textbf{SambaFlow}, integrates with PyTorch and TensorFlow, automatically mapping models to hardware. A Dataflow Graph Analyzer identifies dependencies, while the compiler optimizes execution through parallelization and pipelining. RDA also supports concurrent workloads, allowing simultaneous preprocessing, training, and inference tasks.

The first generation, the \textbf{Cardinal SN10-8R}, featured eight RDUs, each with 640 PCUs and PMUs, and up to 12~TB of configurable DRAM. Its successor, \textbf{Cardinal SN30-8R}, increased compute density to 1,284 PCUs and PMUs per RDU, with 640~MB on-chip memory and 1~TB of DRAM per device. The 2024 \textbf{SN40L-16}~\cite{Prabhakar2024}, optimized for LLM inference, integrates 16 RDUs with 1,040 PCUs and PMUs each, 520~MB of SRAM, and hybrid memory combining 64~GB HBM with 768~GB DDR—significantly boosting bandwidth for compute-bound workloads.


\subsection{Cerebras WSE}
\label{chap:background:sec:cerebras}

Founded in 2015, \textbf{Cerebras Systems} develops large-scale AI accelerators for DNN training and inference. The company has released three systems: the \textbf{CS-1} (2019) with \textbf{WSE-1}, the \textbf{CS-2} (2021) with \textbf{WSE-2}, and the \textbf{CS-3} (2024) featuring \textbf{WSE-3}.

The \textbf{Wafer-Scale Engine (WSE)} is the first processor built on a full 300~mm silicon wafer (462.25~cm$^2$), hosting 400{,}000 cores in WSE-1, 850{,}000 in WSE-2, and 900{,}000 in WSE-3. Each core is a general-purpose compute element supporting arithmetic, logic, and tensor operations through hardware state machines and \textbf{Data Structure Registers (DSRs)} that manage tensor metadata. The 64-bit datapath includes four fused FPMAC units operating on FP16 and bfloat16 data for high tensor throughput.

The WSE adopts a fully distributed memory architecture to maximize locality and bandwidth. Each core integrates 48~KB of SRAM, yielding 44~GB of total on-chip memory in WSE-3. Cores communicate through a 2D mesh of routers, allowing weights and activations to stream efficiently across the wafer and to external \textbf{MemoryX} units, removing the need for shared DRAM or intermediate buffers.

Computation follows a dataflow execution model where operations are triggered by data arrival rather than instruction sequencing. This enables fine-grained sparsity filtering and high energy efficiency. Each core supports dynamic scheduling with up to eight concurrent micro-threads, selecting ready operations based on data availability and priority to improve utilization and reduce latency.


\subsection{Graphcore IPU}
The \textbf{Intelligent Processing Unit (IPU)} by Graphcore is a dataflow-based accelerator designed for AI workloads. Two generations exist: the \textbf{Colossus MK1} (2017) with GC2 IPUs and the \textbf{Colossus MK2} (2020) with GC200 IPUs.

The IPU implements a true \textbf{MIMD} architecture following the \textbf{Bulk Synchronous Parallel (BSP)} model~\cite{Jia2019}, which structures computation into three stages: local execution on private memory, inter-tile communication, and global synchronization. Computation occurs on independent \textbf{tiles}, each with its own processor, local SRAM, and an \textbf{Accumulating Matrix Product (AMP)} unit optimized for matrix multiplications and convolutions. Tiles support up to six threads via SMT, with 1,216 cores (7,296 threads) in the first generation and 1,472 cores (8,832 threads) in the second. Each AMP unit performs up to 16~FP32 or 64~FP16 operations per cycle.

Tiles operate exclusively on local memory—256~KB in the GC2 and 624~KB in the GC200—eliminating caches and shared DRAM while ensuring deterministic performance. Inter-tile communication is handled by the on-chip \textbf{IPU Exchange} and off-chip \textbf{IPU Links}, enabling large-scale, multi-IPU systems.

Programming is performed using the \textbf{Poplar SDK}, where models are represented as computation graphs. The compiler maps graph vertices to threads and edges to communication paths, automatically managing the compute, communication, and synchronization phases under the BSP model.

\begin{table}[h!]
\centering
\footnotesize
\caption{Datacenter-level dataflow AI accelerators from SambaNova, Cerebras, and Graphcore.}
\label{tab:dflow}

\begin{tabularx}{\textwidth}{l l c c c c c c}
\toprule
\textbf{Vendor} & \textbf{Name} & \textbf{Year} & \textbf{Cores} & 
\textbf{TFLOPS} & \textbf{On-chip} & \textbf{Off-chip} & \textbf{TDP} \\
\midrule

\textbf{Graphcore}
& GC2   & 2017 & 1,216 & FP16: 125 
& 300 MB 
& -- 
& 150 W \\

\textbf{Graphcore}
& GC200 & 2020 & 1,472 & FP16: 250 
& 900 MB 
& -- 
& 300 W \\

\midrule

\textbf{SambaNova}
& SN10  & 2021 & 640 
& BF16: N/A 
& 300 MB 
& DRAM: 1.5 TB 
& N/A \\

\textbf{SambaNova}
& SN30  & 2022 & 1,284 
& BF16: 660 
& 640 MB 
& DRAM: 1 TB 
& N/A \\

\textbf{SambaNova}
& SN40L & 2023 & 1,040 
& BF16: 638 
& 520 MB 
& \makecell[l]{HBM: 64 GB \\ DRAM: 768 GB} 
& 900 W \\

\midrule

\textbf{Cerebras}
& CS2 (WSE2) & 2021 & 850K 
& BF16: N/A 
& 41 GB 
& MemoryX 
& 23 kW \\

\textbf{Cerebras}
& CS3 (WSE3) & 2024 & 900K 
& BF16: 125K 
& 44 GB 
& MemoryX 
& 23 kW \\

\bottomrule
\end{tabularx}
\end{table}